\section{Dominio, magnitudes y alcance lógico del certificado computacional de las derivaciones intrínsecas}

\begin{longtable}{@{}>{\raggedright\arraybackslash}p{0.25\textwidth}>{\raggedright\arraybackslash}p{0.65\textwidth}@{}}
\toprule
Bloque & Controles ejecutados\\
\midrule
\endhead
Matriz & 58 filas, 15 campos, identificadores únicos y alcance declarado\\
Vac\'io & orden de \(q_\pm\), positividad, autovalores, cuatro caracteres, identidades constitutivas y completaci\'on modo 30\\
Electrod\'ebil & \(D_A=q_+q_-\) y \(\sin^2\theta_W=1-D_A\)\\
Caos & suma cuadr\'atica, \(899=30^2-1=29\cdot31\), coeficiente cu\'artico, perfil cerrado y residuos num\'ericos\\
Gauss & L\'evy, entrop\'ia y Lochs para bases 10 y 729\\
Vacancias & ventana exhaustiva de \(10^8\) posiciones, huecos \(21/22\), retornos \(6/7\), intervalo certificado y huellas del verificador multiprecisi\'on\\
Espectro & car\'acter \(\chi_{12}\), \(L(1)\), \(L(2)\), unidad hiperb\'olica y torre Bendersky\\
Hamiltoniano modular & \(\Dmod\), pesos 90/120, entrop\'ia, defecto orientado, inversi\'on \`area--memoria, TFD y primera ley\\
Tipo III & grupo de razones \(30\Dmod\Z\) y \(\lambda_\partial=\exp(-30\Dmod)\) bajo hip\'otesis declaradas\\
CKM & inversa de coordenadas, proyector impar, reflexi\'on racional, determinante e intertwining\\
Cosmolog\'ia & eliminaci\'on de \(n\) y ley exacta de potencia seis\\
\bottomrule
\end{longtable}

\section{Valores numéricos certificados de los invariantes constitutivos, espectrales y modulares}

\begin{longtable}{@{}>{\raggedright\arraybackslash}p{0.38\textwidth}>{\raggedright\arraybackslash}p{0.52\textwidth}@{}}
\toprule
Objeto & Valor certificado\\
\midrule
\endhead
\(\widehat\mu\) & \(0.999448350479326935089981555906075895\ldots\)\\
\(\widehat\varepsilon\) & \(0.999999257096636702760441615542323693\ldots\)\\
\(\widehat Z\) & \(0.999724508538937215455554346613611395\ldots\)\\
\(\widehat c\) & \(1.000276310486157526106386268929576268\ldots\)\\
\(D_A\) & \(0.775129119247156430188884096480202784\ldots\)\\
\(1-D_A\) & \(0.224870880752843569811115903519797216\ldots\)\\
\(L(1,\chi_{12})\) & \(0.760345996300946347531094254880405824\ldots\)\\
\(L(2,\chi_{12})\) & \(0.949703126294009395263498491745741516\ldots\)\\
\(\gamma_{\mathbb Q(\sqrt3)}\) & \(1.0539656082484825800\ldots\)\\
\(L\) de L\'evy & \(1.186569110415625452821722975947237121\ldots\)\\
\(h_G\) & \(2.373138220831250905643445951894474241\ldots\)\\
\(I_{\rm or}\) & \(0.001669183233868940655631225912996411\ldots\) nats\\
\(\lambda_\partial\) & \(0.996669646468957378610829976885164731\ldots\)\\
\bottomrule
\end{longtable}

\section{Obstrucciones de tipado y separación de equivalencias entre dominios certificados}

El certificado no promueve las siguientes afirmaciones:

\begin{enumerate}
\item \(I_{\rm or}=\gammaH a_0\); los valores son pr\'oximos pero distintos.
\item El operador del vac\'io es id\'entico al renormalizador de Feigenbaum.
\item Cofinalidad de cilindros implica conjugaci\'on TPK--Gauss.
\item El cierre \(III_{\lambda_\partial}\) vale sin persistencia de ambos canales.
\item La inversi\'on de hoja CKM es la conjugaci\'on CP f\'isica.
\item La realizaci\'on \(\varphi\) y la cronolog\'ia c\'iclica de \(108\) estados
representan la misma cosmolog\'ia.
\item La relaci\'on angular en el esquema de masas f\'isicas incluye por s\'i sola todas las correcciones electrod\'ebiles.
\end{enumerate}

Estos controles negativos no debilitan el volumen: fijan exactamente qu\'e ha sido demostrado y qu\'e mapa adicional cambiar\'ia el estatuto.
